En este capítulo se recogen los resultados de la validación hardware del sistema completo.
Las pruebas se realizaron con la placa CDHS y la placa AOCS conectadas a la ZCU102
mediante FMC. La placa de comunicación serie de diseño propio queda pendiente de
fabricación y sus resultados se incorporarán antes de la entrega final.

\section{Montaje del sistema}

\begin{figure}[H]
    \centering
    \fbox{\parbox{0.85\textwidth}{\centering \textit{[Figura pendiente: foto del sistema completo montado --- ZCU102 con placa CDHS en FMC J5]}\vspace{0.5cm}}}
    \caption{Sistema completo montado: ZCU102 con la placa CDHS conectada al FMC.}
    \label{fig:sistema_montado}
\end{figure}

\begin{figure}[H]
    \centering
    \fbox{\parbox{0.85\textwidth}{\centering \textit{[Figura pendiente: foto de la placa CDHS soldada --- vista superior]}\vspace{0.5cm}}}
    \caption{Placa CDHS con componentes montados (vista superior).}
    \label{fig:cdhs_foto}
\end{figure}

\begin{figure}[H]
    \centering
    \fbox{\parbox{0.85\textwidth}{\centering \textit{[Figura pendiente: foto de la placa AOCS soldada --- vista superior]}\vspace{0.5cm}}}
    \caption{Placa AOCS con componentes montados (vista superior).}
    \label{fig:aocs_foto}
\end{figure}

\section{Validación de comunicaciones serie RS422/RS485}

Las pruebas de comunicaciones serie se realizaron con la aplicación de testing en RTEMS,
que expone una consola interactiva por USB donde el usuario puede enviar mensajes a
cualquier transceptor con el formato \texttt{<ID> <MENSAJE>} o a todos simultáneamente con
\texttt{ALL <MENSAJE>}. Se construyeron arneses de \textit{loopback} para interconectar los canales:

\begin{itemize}
    \item \textbf{RS422}: arnés a 4 hilos cruzando TX\_P/N de un canal con RX\_P/N del opuesto.
    \item \textbf{RS485}: arnés a 2 hilos conectando A+ y B$-$ de un canal con A+ y B$-$ del otro.
\end{itemize}

\subsection{CDHS --- 3 transceptores}

Al arrancar, el driver descubrió automáticamente los 3 transceptores instanciados en el
\textit{bitstream} e imprimió sus direcciones base (\texttt{0xA0000000}, \texttt{0xA0001000},
\texttt{0xA0002000}), confirmando el mecanismo de descubrimiento dinámico de hardware. A
continuación, la consola reportó \texttt{UART 00 [OK]}, \texttt{UART 01 [OK]},
\texttt{UART 02 [OK]}.

Las pruebas de \textit{loopback} mostraron que los mensajes enviados por UART 0 llegaban a
UART 2 y viceversa, y los enviados por UART 1 llegaban correctamente. En las primeras
recepciones de esta captura aparecen caracteres corruptos (\texttt{?}), correspondientes a
una iteración previa del diseño antes de aplicar la corrección de robustez del receptor
descrita en el Capítulo~3 (verificación de \textit{start-bit} a mitad de período). Una
vez implementada esa mejora en la FSM del receptor, los caracteres corruptos desaparecieron
por completo en las pruebas posteriores.

\begin{lstlisting}[caption={Salida de terminal --- prueba RS422/RS485 con la placa CDHS (3 transceptores)}]
[TRANSCEIVER DEBUG] Detectados: 3 | Base: 0xA0000000 | INT: 0xA0003000
UART 00 [OK]  UART 01 [OK]  UART 02 [OK]
CMD> 0 HOLA MUNDO
Tx -> UART 0: OK
[RX UART 02]: HOLA MUNDO
CMD> 2 HOLA MUNDO
Tx -> UART 2: OK
[RX UART 00]: HOLA MUNDO
\end{lstlisting}

\begin{figure}[H]
    \centering
    \fbox{\parbox{0.85\textwidth}{\centering \textit{[Figura pendiente: captura completa del terminal \texttt{cdhs\_testing\_rs} disponible en \texttt{tfm/terminal/}]}\vspace{0.5cm}}}
    \caption{Captura completa del terminal de pruebas RS en la placa CDHS.}
    \label{fig:cdhs_rs_terminal}
\end{figure}

\subsection{AOCS --- 5 transceptores}

Con el \textit{bitstream} de 5 transceptores para la placa AOCS, el descubrimiento detectó
correctamente 5 instancias (\texttt{0xA0000000}--\texttt{0xA0004000}, INTC en
\texttt{0xA0005000}). Las pruebas de \textit{loopback} entre distintos pares de canales
---UART 0$\leftrightarrow$4, UART 0$\leftrightarrow$3, UART 0$\leftrightarrow$2,
UART 0$\leftrightarrow$1--- resultaron todas correctas sin errores de \textit{framing}.

\begin{lstlisting}[caption={Salida de terminal --- prueba RS422/RS485 con la placa AOCS (5 transceptores)}]
[TRANSCEIVER DEBUG] Detectados: 5 | Base: 0xA0000000 | INT: 0xA0005000
CMD> 0 HOLA  ->  [RX UART 04]: HOLA
CMD> 4 HOLA  ->  [RX UART 00]: HOLA
CMD> 0 HOLA  ->  [RX UART 03]: HOLA
CMD> 0 HOLA  ->  [RX UART 01]: HOLA
\end{lstlisting}

\begin{figure}[H]
    \centering
    \fbox{\parbox{0.85\textwidth}{\centering \textit{[Figura pendiente: captura completa del terminal \texttt{aocs\_testing\_rs} disponible en \texttt{tfm/terminal/}]}\vspace{0.5cm}}}
    \caption{Captura completa del terminal de pruebas RS en la placa AOCS.}
    \label{fig:aocs_rs_terminal}
\end{figure}

\subsection{Prueba con 13 transceptores simultáneos}

Para verificar el correcto funcionamiento del driver con un número elevado de instancias, se
cargó un \textit{bitstream} con 13 transceptores instanciados, conectando los canales adicionales a
los pines externos del conector J3 de la ZCU102. El sistema arrancó y operó correctamente
con todos los canales, confirmando que la arquitectura de ISR maestra con tareas \textit{worker}
individuales escala sin problemas hasta el número máximo de instancias soportado. La
prueba con los 14 transceptores sobre la placa de comunicación serie de diseño propio queda
pendiente de fabricación.

\section{Validación del bus CAN}

Las pruebas del bus CAN se realizaron con la aplicación de test desarrollada por Diego Ramos
(compañero del laboratorio B105 en el proyecto LINCE), que configura los canales CAN y
ejecuta una batería de tests de enlace.

La suite de tests ejecutó 26 pruebas cubriendo inicialización, \textit{loopback} interno,
transferencia física CAN0$\leftrightarrow$CAN1, filtrado de identificadores estándar y
extendido, manejo de tramas RTR y pruebas de interrupción por hardware. El resultado fue
\textbf{26/26 tests PASS, 0 FAILED}:

\begin{lstlisting}[caption={Resumen de resultados del test CAN (placa CDHS)}]
[PASS] CAN0 initialization (Fast Mode)
[PASS] CAN1 initialization (Slow Mode)
[PASS] Loopback RX ID matches TX ID
[PASS] Loopback RX Data matches TX Data
[PASS] CAN0 received correct ID from CAN1        (fisico)
[PASS] CAN0 received correct Data from CAN1      (fisico)
[PASS] Standard ID: Exact Match Accepted (0x1A4)
[PASS] Extended ID: Exact Match Accepted (0x12345678)
[PASS] RTR Filter: Accepted matching Remote Request
[PASS] TxOk fired successfully
[PASS] RxOk fired and FIFO was drained successfully (No Storm)
... (26/26 PASS, 0 FAILED)
\end{lstlisting}

Un resultado relevante adicional fue la verificación del efecto de las resistencias de
terminación sobre la integridad de la señal CAN. Sin terminación, la señal presentaba
reflexiones visibles que degradaban los flancos; con ambas resistencias de 60~$\Omega$
conectadas, la forma de onda resultó limpia y bien definida.

\begin{figure}[H]
    \centering
    \fbox{\parbox{0.85\textwidth}{\centering \textit{[Figura pendiente: captura de osciloscopio --- señal CAN sin resistencias de terminación]}\vspace{0.5cm}}}
    \caption{Señal diferencial CAN sin resistencias de terminación.}
    \label{fig:can_sin_term}
\end{figure}

\begin{figure}[H]
    \centering
    \fbox{\parbox{0.85\textwidth}{\centering \textit{[Figura pendiente: captura de osciloscopio --- señal CAN con resistencias de terminación]}\vspace{0.5cm}}}
    \caption{Señal diferencial CAN con resistencias de terminación (60~$\Omega$ + 60~$\Omega$).}
    \label{fig:can_con_term}
\end{figure}

\section{Validación del ADC SPI (termistores CDHS)}

La aplicación de lectura del ADC ADS7950 muestreó los cuatro canales a 500~ms de intervalo
e imprimió los valores digitales por el terminal. Para verificar la linealidad de la cadena de
adquisición se aplicaron tensiones de referencia conocidas a las entradas analógicas:

\begin{itemize}
    \item Tensión de entrada 0~V $\rightarrow$ valor digital próximo a 0 (fondo de escala inferior).
    \item Tensión de entrada 3,3~V $\rightarrow$ valor digital próximo a 4095 (fondo de escala
    superior, 12 bits).
\end{itemize}

Los resultados confirmaron el rango de conversión esperado, validando el driver SPI de bajo
nivel y la cadena analógica de la placa CDHS. La captura registra el canal CH2 siendo
sometido a un barrido de tensión con la fuente de laboratorio:

\begin{lstlisting}[caption={Lectura del ADC ADS7950 durante barrido de tensión en CH2}]
ADS7950: CH0= 244  CH1=  43  CH2=   0  CH3=  19   <- tension en CH2 = 0 V
ADS7950: CH0= 245  CH1=  43  CH2=1577  CH3=  19   <- subida
ADS7950: CH0= 245  CH1=  43  CH2=3419  CH3=  21   <- cerca del maximo
ADS7950: CH0= 245  CH1=  43  CH2=3565  CH3=  21   <- ~2.9 V aplicados
ADS7950: CH0= 245  CH1=  43  CH2=2436  CH3=  22   <- bajada
ADS7950: CH0= 246  CH1=  43  CH2= 102  CH3=  22   <- de vuelta a 0 V
\end{lstlisting}

Los canales CH0 ($\approx$245), CH1 ($\approx$43) y CH3 ($\approx$19) mostraron valores
estables bajos durante todo el ensayo, correspondientes a las entradas no excitadas.

\begin{figure}[H]
    \centering
    \fbox{\parbox{0.85\textwidth}{\centering \textit{[Figura pendiente: captura completa del terminal \texttt{cdhs\_testing\_adc.txt} disponible en \texttt{tfm/terminal/}]}\vspace{0.5cm}}}
    \caption{Captura completa del terminal de lectura ADC (placa CDHS).}
    \label{fig:cdhs_adc_terminal}
\end{figure}

\section{Validación PWM (calentadores CDHS y puentes en H AOCS)}

Las señales PWM generadas por el bloque \texttt{PWMx4\_auto\_test} se midieron con el
osciloscopio sobre el conector J5 de la placa CDHS, verificando los cuatro canales: 10~kHz,
5~kHz, 1~kHz y 100~Hz, todos con \textit{duty cycle} del 50\%.

Para la placa AOCS, el bloque VHDL de control de motores generó señales PWM alternando
la dirección de giro cada pocos segundos (línea 1 activa con línea 2 a 0, y a continuación
línea 1 a 0 con línea 2 activa), permitiendo verificar el control de puentes en H con una
fuente de alimentación externa conectada a los bornes banana.

\begin{figure}[H]
    \centering
    \fbox{\parbox{0.85\textwidth}{\centering \textit{[Figura pendiente: captura de osciloscopio --- señal PWM en J5 de la placa CDHS (frecuencia y duty cycle)]}\vspace{0.5cm}}}
    \caption{Señal PWM medida en el conector J5 de la placa CDHS.}
    \label{fig:cdhs_pwm}
\end{figure}

\begin{figure}[H]
    \centering
    \fbox{\parbox{0.85\textwidth}{\centering \textit{[Figura pendiente: captura de osciloscopio --- par de señales PWM complementarias de control de motor (AOCS)]}\vspace{0.5cm}}}
    \caption{Señales PWM de control de motor en la placa AOCS (par complementario para puente en H).}
    \label{fig:aocs_pwm}
\end{figure}
